% 埃米·诺特（综述）
% license CCBYSA3
% type Wiki

本文根据 CC-BY-SA 协议转载翻译自维基百科\href{https://en.wikipedia.org/wiki/Emmy_Noether}{相关文章}

\begin{figure}[ht]
\centering
\includegraphics[width=6cm]{./figures/c0d48f5d45fa6e41.png}
\caption{诺特，约摄于1900年至1910年间。} \label{fig_AMnt_1}
\end{figure}
阿玛莉·艾米·诺特（Amalie Emmy Noether，\(^\text{[a][b]}\)，1882年3月23日－1935年4月14日）是一位德国数学家，她在抽象代数领域作出了许多重要贡献。她还证明了诺特第一与第二定理，这两项成果在数学物理中具有基础性意义。\(^\text{[4]}\)帕维尔·亚历山德罗夫、阿尔伯特·爱因斯坦、让·迪厄多内、赫尔曼·外尔以及诺伯特·维纳都曾称她为“数学史上最重要的女性”。\(^\text{[5][6][7]}\)作为当时最杰出的数学家之一，诺特发展了环论、域论和代数理论。在物理学中，诺特定理揭示了对称性与守恒定律之间的深刻联系。\(^\text{[8]}\)

诺特出生于法兰克尼亚小镇埃尔朗根的一个犹太家庭，她的父亲是数学家马克斯·诺特。起初，她计划在通过所需考试后教授法语和英语，但后来转而在其父任教的埃尔朗根－纽伦堡大学学习数学。1907年，在保罗·戈尔丹的指导下，她完成了博士论文。此后，她在埃尔朗根数学研究所无薪工作了七年。\(^\text{[9]}\)当时，女性基本上被排除在学术职位之外。1915年，大卫·希尔伯特和费利克斯·克莱因邀请她加入哥廷根大学数学系——这是当时世界知名的数学研究中心。然而哲学学院对此提出反对，因此她在接下来的四年里只能以希尔伯特的名义讲课。1919年，她的任教资格最终获批，从而获得了“私人讲师”的职称。\(^\text{[9]}\)

诺特一直是哥廷根大学数学系的核心成员，直至1933年为止；她的学生们有时被称为“诺特男孩”。1924年，荷兰数学家B.L.范德瓦尔登加入她的学术圈，并很快成为她思想的主要阐述者；她的研究成果成为他1931年具有深远影响的教材《现代代数》第二卷的理论基础。到1932年在苏黎世举行的国际数学家大会上作大会报告时，诺特的代数才能已获得世界范围的认可。次年，德国纳粹政府将犹太人清除出大学职位，诺特被迫移居美国，在宾夕法尼亚州的布林茅尔学院任教。在那里，她教授研究生和博士后课程，学生包括玛丽·约翰娜·魏斯和奥尔加·陶斯基－托德。与此同时，她还在新泽西州普林斯顿的高等研究院讲学并从事研究工作。\(^\text{[9]}\)

诺特的数学研究通常被划分为三个“时期”。\(^\text{[10]}\)在第一个时期（1908–1919年），她对代数不变量理论和数域理论作出了重要贡献。她在变分法中的微分不变量研究——即著名的诺特定理——被称为“在指导现代物理学发展过程中最重要的数学定理之一”。\(^\text{[11]}\)在第二个时期（1920–1926年），她开始从事一项“改变了[抽象]代数面貌”的工作。\(^\text{[12]}\)在她1921年的经典论文《环域中的理想论》（Idealtheorie in Ringbereichen，意为“环域中理想的理论”）中，诺特将交换环中的理想理论发展成为一个具有广泛应用的数学工具。她巧妙地运用了升链条件，而满足该条件的数学对象也因纪念她而被称为“诺特环”。在第三个时期（1927–1935年），她发表了关于非交换代数与超复数的研究成果，并将群的表示论与模论和理想论统一起来。除了自己的论文之外，诺特还慷慨地分享思想，对多项由其他数学家发表的研究方向作出了重要贡献，甚至包括一些与她主要研究领域相距甚远的领域，如代数拓扑。
\subsection{传记}
\subsubsection{早年生活}
\begin{figure}[ht]
\centering
\includegraphics[width=6cm]{./figures/9bf717f586e25b88.png}
\caption{诺特在巴伐利亚的埃尔朗根市长大，图为1916年的一张该城市明信片。} \label{fig_AMnt_2}
\end{figure}
阿玛莉·艾米·诺特于1882年3月23日出生在巴伐利亚的埃尔朗根。\(^\text{[13]}\)她是数学家马克斯·诺特与伊达·阿玛莉亚·考夫曼的四个孩子中的长女，父母皆出身于富裕的犹太商人家庭。\(^\text{[14]}\)她的名字原本是“阿玛莉”，但从小便开始使用中间名“艾米”，并在成年后及其所有学术著作中始终沿用这一名字。\(^\text{[b]}\)

年轻时的诺特在学业上并不特别突出，但以聪慧与友善著称。她有些近视，童年时期说话带轻微的口齿不清。多年后，一位家庭朋友回忆说，在一次儿童聚会上，小诺特曾迅速解开一道益智谜题，显示出她早期出众的逻辑才能。\(^\text{[15]}\)她像当时大多数女孩一样学习烹饪与家务，也上过钢琴课。虽然她对这些活动都缺乏热情，但却非常热爱舞蹈。\(^\text{[16]}\)
\begin{figure}[ht]
\centering
\includegraphics[width=6cm]{./figures/40e1cdf2d55a10d2.png}
\caption{艾米·诺特与她的三位弟弟——阿尔弗雷德、弗里茨和罗伯特，摄于1918年之前。} \label{fig_AMnt_3}
\end{figure}
诺特有三个弟弟。长弟阿尔弗雷德·诺特生于1883年，1909年在埃尔朗根获得化学博士学位，但九年后便去世。\(^\text{[17]}\)次弟弗里茨·诺特生于1884年，在慕尼黑求学，并在应用数学领域作出过贡献。\(^\text{[18]}\)据推测，他于第二次世界大战期间的1941年在苏联遭到处决。\(^\text{[19]}\)最小的弟弟古斯塔夫·罗伯特·诺特生于1889年，关于他的生平所知甚少；他长期患有慢性疾病，并于1928年去世。\(^\text{[20][21]}\)
\subsubsection{教育经历}
诺特自幼在法语与英语方面展现出优秀的天赋。1900年初，她参加了外语教师资格考试，并获得了“Sehr gut”（非常好）的总评成绩。这一成绩使她有资格在女子学校教授外语，但她选择继续在父亲任教的埃尔朗根－纽伦堡大学深造数学。[22][23]

这一决定在当时极不寻常。就在两年前，该校学术参议院曾宣称，允许男女同校将会“颠覆整个学术秩序”。\(^\text{[24]}\)在这所拥有986名学生的大学里，诺特是仅有的两位女性之一。她只能旁听课程，不能正式注册听课，并且每听一门课都必须得到任课教授的特别许可。尽管困难重重，她仍于1903年7月14日通过了纽伦堡实科中学的毕业考试。\(^\text{[22][25][26]}\)

在1903至1904年的冬季学期中，她前往哥廷根大学学习，听取了天文学家卡尔·施瓦西以及数学家赫尔曼·闵可夫斯基、奥托·布卢门塔尔、费利克斯·克莱因和大卫·希尔伯特的讲座。\(^\text{[27]}\)
\begin{figure}[ht]
\centering
\includegraphics[width=6cm]{./figures/8f7a5161e1329585.png}
\caption{保罗·戈尔丹指导了诺特的博士论文，研究主题为四次型不变量。} \label{fig_AMnt_4}
\end{figure}
1903年，巴伐利亚大学对女性全面入学的限制被取消。\(^\text{[28]}\)诺特返回埃尔朗根，并于1904年10月正式重新入学，声明她将专注于数学研究。当年她所在年级共有六名女性（其中两名为旁听生），而她是唯一一位选择数学专业的女性。\(^\text{[29]}\)在保罗·戈尔丹的指导下，她于1907年完成了博士论文《三元四次型不变量系的构成》，并于同年以“最优等成绩”毕业。\(^\text{[30][31]}\)戈尔丹属于“不变量计算学派”的代表人物，而诺特的论文以列出三百多个明确计算出的不变量作为结尾。这种具体计算的方法后来被希尔伯特所开创的更抽象、更一般化的不变量理论方法所取代。\(^\text{[32][33]}\)尽管她的论文在当时受到好评，但诺特后来却将这篇论文以及她之后撰写的几篇类似论文称为“垃圾”。\(^\text{[33]}\)她此后的全部研究工作都进入了一个完全不同的领域。\(^\text{[34]}\)
\subsubsection{埃尔朗根－纽伦堡大学}
1908年至1915年间，诺特在埃尔朗根大学数学研究所任教，但没有任何薪酬；在其父马克斯·诺特身体欠佳无法授课时，她偶尔会代替父亲讲课。\(^\text{[35]}\)她于1908年加入巴勒莫数学学会，并于1909年加入德国数学家协会。\(^\text{[36]}\)在1910年和1911年间，她发表了将自己论文研究内容从三元变量推广至$n$元变量的成果。\(^\text{[37]}\)
\begin{figure}[ht]
\centering
\includegraphics[width=6cm]{./figures/7f692954ec96cc8c.png}
\caption{诺特有时会用明信片与同事恩斯特·菲舍尔讨论抽象代数问题。此明信片的邮戳日期为1915年4月10日。} \label{fig_AMnt_5}
\end{figure}
戈尔丹于1910年退休，\(^\text{[38]}\)随后诺特在他的继任者——埃尔哈德·施密特与恩斯特·菲舍尔手下任教，后者于1911年接替了前者的职位。\(^\text{[39]}\)据她的同事赫尔曼·外尔及其传记作者奥古斯特·迪克记载，菲舍尔对诺特产生了重要影响，尤其是他将大卫·希尔伯特的研究成果介绍给了她。\(^\text{[40][41]}\)诺特与菲舍尔都对数学充满热情，他们常常在讲座结束后仍意犹未尽地讨论问题。据悉，诺特甚至会给菲舍尔寄去明信片，继续推演她的数学思路。\(^\text{[42][43]}\)

1913年至1916年间，诺特发表了多篇论文，将希尔伯特的方法加以推广与应用，研究对象包括有理函数域以及有限群的不变量。\(^\text{[44]}\)这一阶段标志着她首次涉足抽象代数——一个她后来将作出开创性贡献的领域。\(^\text{[45]}\)

在埃尔朗根期间，诺特指导了两名博士生：\(^\text{[46]}\)汉斯·法尔肯贝格与弗里茨·塞德尔曼，他们分别于1911年与1916年完成论文答辩。\(^\text{[47][48]}\)尽管诺特在指导中发挥了重要作用，但两人的官方导师名义上仍是她的父亲。博士毕业后，法尔肯贝格先后在不伦瑞克与柯尼斯堡任职，后成为吉森大学教授，\(^\text{[49]}\)而塞德尔曼则在慕尼黑任教。\(^\text{[46]}\)
\subsubsection{哥廷根大学}
\textbf{资格授予与诺特定理}

1915年初，大卫·希尔伯特与费利克斯·克莱因邀请诺特回到哥廷根大学任教。然而，他们的聘请计划最初遭到哲学学院中语言学家与历史学家的反对，他们坚持认为女性不应被任命为“私人讲师”。在一次关于此事的联合系会议上，一位教授抗议道：“我们的士兵回到大学后，如果发现自己必须在一位女性脚下学习，他们会作何感想？”\(^\text{[50][51]}\)希尔伯特对此极为愤慨，他认为唯一重要的标准是学术资格，而与性别无关。他斥责那些反对授予诺特任教资格的人。虽然他当时的原话未被完整记录，但人们普遍认为他在辩论中讽刺道，大学“不是澡堂”。\(^\text{[40][50][52][53]}\)根据帕维尔·亚历山德罗夫的回忆，学院成员对诺特的反对不仅源于性别歧视，也与他们对她的社会民主主义政治立场及犹太血统的偏见有关。\(^\text{[53]}\)
\begin{figure}[ht]
\centering
\includegraphics[width=6cm]{./figures/ae8a0eff4b92d74a.png}
\caption{大卫·希尔伯特于1915年邀请诺特加入哥廷根大学数学系，此举挑战了部分同事的观念——他们认为女性不应在大学任教。} \label{fig_AMnt_6}
\end{figure}
诺特于4月下旬前往哥廷根；两周后，她的母亲在埃尔朗根突然去世。此前，她母亲曾因眼疾接受治疗，但这种疾病的具体性质及其与去世的关联仍不明。大约在同一时期，诺特的父亲退休，而她的弟弟参军加入德军，参加第一次世界大战。诺特返回埃尔朗根待了数周，主要是为了照顾年迈的父亲。\(^\text{[54]}\)

在她最初在哥廷根任教的几年里，诺特没有正式职位，也没有薪水。她的讲座常常以希尔伯特的名义对外发布，而她自己仅被列为“助讲”。\(^\text{[55]}\)

到达哥廷根后不久，她便以一项非凡的成果展现了自己的才能——她证明了如今被称为诺特定理的命题：任何可微对称性都对应着一个物理系统的守恒定律。\(^\text{[51][56]}\)她的论文《不变变分问题》由同事费利克斯·克莱因代为宣读，于1918年7月26日在哥廷根皇家科学院会议上发表。\(^\text{[57][58]}\)诺特本人当时可能因并非学会成员而无法亲自宣读。\(^\text{[59]}\)美国物理学家莱昂·莱德曼与克里斯托弗·T·希尔在其著作《对称与优美的宇宙》中指出，诺特定理“无疑是指导现代物理学发展的最重要数学定理之一，其重要性或可与毕达哥拉斯定理相提并论”。\(^\text{[11]}\)
\begin{figure}[ht]
\centering
\includegraphics[width=8cm]{./figures/2d2fd4171cdfdeb2.png}
\caption{哥廷根大学于1919年批准了诺特的任教资格申请，此时距离她在该校开始讲学已过去四年。} \label{fig_AMnt_7}
\end{figure}
第一次世界大战结束后，1918至1919年的德国革命带来了社会观念的显著变化，其中包括女性权利的提升。1919年，哥廷根大学允许诺特正式申请取得“任教资格”（habilitation，相当于终身教职的资格）。她的口试于当年5月底举行，并于6月成功完成资格讲座。\(^\text{[60]}\)由此，诺特成为一名私人讲师，\(^\text{[61]}\)同年秋季学期，她首次以自己的名字开设课程。\(^\text{[62]}\)尽管如此，她的教学工作仍未获得薪酬。\(^\text{[55]}\)

三年后，普鲁士科学、艺术与公共教育部长奥托·博利茨（Otto Boelitz [de]）致信诺特，授予她“非公务员编制的特别教授”的头衔。\(^\text{[63]}\)这一称号相当于“未聘终身的副教授”，仅赋予她有限的校内行政职权。这是对她学术成就的重要认可，但依然不包含薪资待遇。直到一年后，她被任命为“代数讲师”的特别职位，才开始正式获得报酬。\(^\text{[64][65]}\)

\textbf{抽象代数研究}

诺特定理对经典力学与量子力学产生了深远影响，但在数学界，她最为人铭记的贡献在于抽象代数领域。纳森·雅各布森在《诺特论文集》序言中写道：\\\\
“抽象代数的发展——二十世纪数学中最具独创性的革新之一——在很大程度上归功于她，无论是在发表的论文、讲座，还是她对同时代人的学术影响中。”\(^\text{[66]}\)

诺特在代数方面的研究始于1920年。当时她与自己的学生兼合作者维尔纳·施迈德勒共同发表了一篇关于理想理论的论文，在其中他们首次定义了环中的左理想与右理想。\(^\text{[45]}\)

次年，她发表了著名论文《环域中的理想论》，\(^\text{[67]}\)系统分析了与数学理想相关的升链条件，并在此基础上完整地证明了拉斯克－诺特定理的普遍形式。著名代数学家欧文·卡普兰斯基称这项工作是“革命性的”。[68] 该论文的发表使得满足升链条件的数学对象被称为诺特对象，以纪念她的贡献。[68][69]
\begin{figure}[ht]
\centering
\includegraphics[width=6cm]{./figures/6c05e0ad0960ebae.png}
\caption{B.L. 范德瓦尔登（摄于1980年）在哥廷根深受诺特的影响。} \label{fig_AMnt_8}
\end{figure}
1924年，一位年轻的荷兰数学家巴特尔·伦德特·范德瓦尔登来到哥廷根大学。他立即开始与诺特合作，诺特向他传授了极具价值的抽象概念化方法。范德瓦尔登后来评价说，她的原创性“绝对无可比拟”。\(^\text{[70]}\)回到阿姆斯特丹后，他撰写了代数学领域的重要教材《现代代数》，该书共两卷，是该领域的奠基性著作；其第二卷于1931年出版，深受诺特研究成果的启发。\(^\text{[71]}\)诺特本人从未主动寻求名誉，但在该书第七版中，范德瓦尔登特意注明：“部分内容基于E. 阿廷与E. 诺特的讲座。”\(^\text{[72][73][74]}\)自1927年起，诺特开始与埃米尔·阿廷、理查德·布劳尔以及赫尔穆特·哈塞密切合作，研究非交换代数问题。\(^\text{[40][71]}\)

范德瓦尔登的到访正值世界各地的数学家纷纷汇聚哥廷根，这里当时已成为数学与物理研究的国际中心。1923年，俄罗斯数学家帕维尔·亚历山德罗夫与帕维尔·乌雷松是最早到来的几位学者。\(^\text{[75]}\)1926年至1930年间，亚历山德罗夫经常在哥廷根大学讲学，并与诺特结下深厚友谊。\(^\text{[76]}\)他亲切地称她为“der Noether”，这里的“der”并非德语阳性冠词，而是一种带敬意的昵称。\(^\text{[c][76]}\)诺特曾努力为亚历山德罗夫争取在哥廷根的常任教授职位，但最终只成功为他从洛克菲勒基金会获得了一笔奖学金，使他得以于1927至1928学年赴普林斯顿大学进修。\(^\text{[76][79]}\)

\textbf{研究生}
\begin{figure}[ht]
\centering
\includegraphics[width=6cm]{./figures/47ccbbdb2f12c29b.png}
\caption{诺特，约摄于1930年左右。} \label{fig_AMnt_9}
\end{figure}
在哥廷根期间，诺特指导了十多名博士生。\(^\text{[46]}\)由于当时规定她不得单独指导论文，大多数学生的论文名义上由埃德蒙·朗道等人共同指导。\(^\text{[80][81]}\)她的第一位博士生是格蕾特·赫尔曼，后者于1925年2月完成论文答辩。\(^\text{[82]}\)虽然赫尔曼后来因其在量子力学基础方面的研究而更为人熟知，但她的博士论文在理想理论领域被认为具有重要贡献。\(^\text{[83][84]}\)赫尔曼后来满怀敬意地称诺特为自己的“论文之母”。\(^\text{[82]}\)

大约在同一时期，海因里希·格雷尔和鲁道夫·黑尔策也在诺特的指导下撰写论文。黑尔策在论文答辩前因肺结核去世。\(^\text{[82][85][86]}\)格雷尔于1926年完成答辩，随后在耶拿大学与哈雷大学任职，但在1935年因被指控有同性恋行为而被吊销教职。\(^\text{[46]}\)后来他获得平反，并于1948年成为洪堡大学教授。\(^\text{[46][82]}\)

随后，诺特指导了维尔纳·韦伯\(^\text{[87]}\)和雅各布·莱维茨基\(^\text{[88]}\)，两人均于1929年完成论文答辩。\(^\text{[89][90]}\)韦伯被认为是一位能力平平的数学家，\(^\text{[80]}\)后来却参与了将犹太数学家驱逐出哥廷根的行动。\(^\text{[91]}\)莱维茨基先后在耶鲁大学与当时由英国托管的巴勒斯坦地区的耶路撒冷希伯来大学任职，他在环论方面作出了重要贡献，尤其以莱维茨基定理和霍普金斯－莱维茨基定理而著称。\(^\text{[90]}\)

其他被称为“诺特男孩”的学生包括马克斯·德林、汉斯·菲廷、恩斯特·维特、曾炯之和奥托·席林。其中，德林被认为是诺特最有前途的学生之一，他于1930年获得博士学位。\(^\text{[92][93]}\) 他先后在汉堡、马尔登和哥廷根工作，\(^\text{[d]}\)并以其在算术几何领域的贡献而闻名。\(^\text{[95]}\)菲廷于1931年以关于阿贝尔群的论文毕业，\(^\text{[96]}\)后因在群论领域的研究成果被人铭记，特别是提出了菲廷定理与菲廷引理。\(^\text{[97]}\)他因骨疾早逝，年仅31岁。\(^\text{[98]}\)

维特最初由诺特指导，但在1933年4月诺特的职位被撤销后，他被改由古斯塔夫·赫尔格洛茨指导。\(^\text{[98]}\)他于1933年7月获得博士学位，论文主题为黎曼－罗赫定理与ζ函数，\(^\text{[99]}\)此后在数学领域作出了多项以他命名的贡献。\(^\text{[97]}\)曾炯之因证明曾定理而被后人铭记，他于当年12月获得博士学位。\(^\text{[100]}\)1935年他返回中国，在国立浙江大学任教，\(^\text{[97]}\)但仅五年后便英年早逝。\(^\text{[e]}\)奥托·席林同样曾师从诺特，但因诺特被迫流亡海外而不得不另寻导师。他后来在赫尔穆特·哈塞的指导下，于1934年在马尔堡大学完成博士学位。\(^\text{[97][102]}\)此后，他在剑桥大学三一学院任博士后研究员，之后移居美国工作。\(^\text{[46]}\)

诺特的其他学生还包括：威廉·德恩特，于1927年以一篇关于群论的论文获得博士学位；\(^\text{[103]}\)维尔纳·福尔贝克，于1935年以分裂域为题完成博士论文；\(^\text{[46]}\)以及沃尔夫冈·维希曼，于1936年以p进理论为主题获得博士学位。\(^\text{[104]}\)关于前两人几乎没有资料，但据悉维希曼曾支持学生发起的一项未成功的请愿，试图撤销对诺特的解聘决定。\(^\text{[105]}\)他后来作为士兵在第二次世界大战的东线战场上阵亡。\(^\text{[46]}\)

\textbf{诺特学派}

除了她的博士生之外，诺特还培养了一批与她志同道合的数学家，他们共同推动了抽象代数的发展。\(^\text{[106]}\)这一群体常被称为“诺特学派”。\(^\text{[107][108]}\)例如，她与沃尔夫冈·克鲁尔的密切合作极具代表性。克鲁尔通过提出主理想定理以及交换环维数理论，极大地推进了交换代数的发展。[109] 另一位重要成员是戈特弗里德·克特，他借助诺特与克鲁尔的方法，在超复数量理论领域作出了重要贡献。\(^\text{[109]}\)

除了卓越的数学洞察力外，诺特也因其待人真诚而受到尊敬。她有时对持不同意见者表现得直率甚至有些无礼，但总体上以乐于助人、耐心指导新生而闻名。她对数学精确性的极高要求使一位同事称她为“严厉的批评者”，但这种严格总是与温和的教导并存。\(^\text{[110]}\)

在范德瓦尔登为她撰写的讣告中，他这样评价诺特：

“她完全没有自我中心与虚荣之心，从不为自己争取任何荣誉，而总是把学生的成果放在首位。”\(^\text{[70]}\)

诺特对学科与学生的投入远超出一般学术范畴。一次，学校因国定假日关闭，她便带着学生们在教学楼外的台阶集合，穿过树林，到附近的咖啡馆继续讲课。\(^\text{[1111]}\)后来，在被纳粹政权解职后，她仍邀请学生到家中讨论他们的未来计划与数学问题。\(^\text{[112]}\)

\textbf{有影响力的讲座}

诺特节俭的生活方式最初源于她长期得不到薪酬。即使在1923年大学开始为她发放少量工资后，她依然过着简单朴素的生活。她在晚年收入虽有所增加，但仍将一半薪水存起来，留给她的侄子戈特弗里德·E·诺特。\(^\text{[113]}\)

传记作者指出，诺特对外表与举止几乎毫不在意，她的全部注意力都集中在学术上。著名代数学家、曾受教于她的奥尔加·陶斯基－托德曾回忆一次午餐场景：诺特当时全神贯注于数学讨论，“双手挥舞着比划”，一边吃饭一边“不断把食物洒出来，再若无其事地从衣服上擦掉”。\(^\text{[114]}\)那些注重仪表的学生在课堂上常感到尴尬——她会随手从衣襟里取出手帕，毫不在意头发在讲课中越发凌乱。一次，在长达两小时的课程间隙，两个女学生试图上前提醒她注意形象，但最终没能插上话，因为她正热烈地与其他学生探讨数学问题。\(^\text{[115]}\)

诺特的讲课从不遵循固定的教学计划。\(^\text{[70]}\)她语速很快，课程内容对许多人而言都相当难以跟上，其中包括卡尔·路德维希·齐格尔和保罗·迪布雷伊。\(^\text{[116][117]}\)不喜欢她讲课风格的学生常感到疏离。\(^\text{[118]}\)那些偶尔旁听她课程的“外来听众”往往在教室里坐不到半小时就因困惑或挫败而离开。一位常听课的学生曾形象地描述道：“敌人被击退了——他逃走了。”\(^\text{[118]}\)

诺特把课堂当作与学生即兴讨论的场所，用以共同思考与澄清数学中的关键问题。她的一些最重要成果正是在这样的课堂讨论中形成的，而学生们整理的课堂笔记后来成为多部重要教材的基础，例如范德瓦尔登与德林的著作。\(^\text{[70]}\)诺特那种具有感染力的数学热情深深影响了最投入的学生，他们珍视与她的生动讨论，将这种思想碰撞视为真正的学术享受。\(^\text{[120][121]}\)

她的多位同事常来旁听她的课程，而诺特有时也会允许他人（包括她的学生）署名发表源自她思想的成果，因此她的许多研究最终以他人名义发表于论文中。\(^\text{[71][72]}\)据记录，诺特在哥廷根大学至少开设过以下五门为期一学期的课程：\(^\text{[122]}\)
\begin{itemize}
\item 1924–1925年冬季学期：Gruppentheorie und hyperkomplexe Zahlen（群论与超复数）
\item 1927–1928年冬季学期：Hyperkomplexe Grössen und Darstellungstheorie（超复数量与表示论）
\item 1928年夏季学期：Nichtkommutative Algebra（非交换代数）
\item 1929年夏季学期：Nichtkommutative Arithmetik（非交换算术）
\item 1929–1930年冬季学期：Algebra der hyperkomplexen Grössen（超复数量代数）
\end{itemize}
\subsubsection{莫斯科国立大学}
\begin{figure}[ht]
\centering
\includegraphics[width=6cm]{./figures/8b2db2722c8ba117.png}
\caption{诺特于1928年至1929年在莫斯科国立大学任教。} \label{fig_AMnt_10}
\end{figure}
1928年至1929年间，诺特接受邀请前往莫斯科国立大学任教，在那里继续与帕维尔·亚历山德罗夫合作。除了继续她的研究工作外，她还开设了抽象代数与代数几何课程。她曾与拓扑学家列夫·庞特里亚金和尼古拉·切博塔廖夫共事，后者称赞她在伽罗瓦理论发展中的重要贡献。\(^\text{[123][124][125]}\)
\begin{figure}[ht]
\centering
\includegraphics[width=6cm]{./figures/1545a6308f8fd4f1.png}
\caption{帕维尔·亚历山德罗夫} \label{fig_AMnt_11}
\end{figure}
虽然政治并非她生活的中心，但诺特对政治事务始终抱有浓厚兴趣。据亚历山德罗夫回忆，她对俄国革命表示出极大支持。她尤其欣喜于苏联在科学与数学领域取得的进步，认为这些成果体现了布尔什维克事业所创造的新机会。这种态度在德国为她带来了一些麻烦，最终导致她被驱逐出一家宿舍楼，因为学生领袖抱怨说他们“与一位倾向马克思主义的犹太女人同住”。\(^\text{[126]}\)赫尔曼·外尔回忆道：“在1918年革命后的动荡时期，诺特或多或少站在社会民主党的立场上。”\(^\text{[401]}\)她于1919年至1922年间是“独立社会民主党”成员——这是一个寿命不长的分裂政党。逻辑学家兼数学史学家科林·麦克拉蒂评价道：“她并非布尔什维克，但也并不畏惧被人那样称呼。”\(^\text{[127]}\)

诺特原本计划重返莫斯科，并得到了亚历山德罗夫的支持。1933年她离开德国后，亚历山德罗夫曾尝试通过苏联教育部为她争取莫斯科国立大学的教授职位。尽管此事未能成功，但两人在1930年代持续通信，而她也在1935年筹划重返苏联的行程。\(^\text{[126]}\)
\subsubsection{荣誉与认可}
1932年，艾米·诺特与埃米尔·阿廷共同获得了阿克曼－托伊布纳纪念奖，以表彰他们在数学领域的杰出贡献。\(^\text{[71]}\)该奖项附带500帝国马克（ℛ︁ℳ︁）的奖金，被视为对她在代数学领域长期卓越工作的迟来认可。然而，她的同事们仍为她未能当选哥廷根科学院成员、也从未被晋升为正式教授（Ordentlicher Professor，即正教授）而感到遗憾与不满。\(^\text{[63][128][129]}\)

1932年，诺特迎来了她的五十岁生日，同事们以典型的数学家方式为她庆祝。赫尔穆特·哈塞在《数学年刊》上发表了一篇献给她的论文，其中通过证明一个非交换互反律，印证了她此前的猜想——某些非交换代数的性质确实比交换代数更为简洁。\(^\text{[130]}\)这一成果令她非常欣喜。哈塞还给她寄去一个数学谜题，他称之为“音节的 mμν 之谜”。诺特很快就解出了它，但可惜这一谜题后来已遗失。\(^\text{[128][129]}\)
\begin{figure}[ht]
\centering
\includegraphics[width=8cm]{./figures/f1e6cdc1213aa6d8.png}
\caption{1932年，诺特前往苏黎世，在国际数学家大会上发表大会报告。} \label{fig_AMnt_12}
\end{figure}
同年9月，诺特在苏黎世举行的国际数学家大会上发表大会报告（德语称 großer Vortrag），题为《超复系统与交换代数及数论的关系》。大会共有800名与会者，其中包括她的同事赫尔曼·外尔、埃德蒙·朗道和沃尔夫冈·克鲁尔。会议共有420位正式参会者，并进行了21场大会报告。诺特获邀作大会报告这一显赫位置，显然是对她在数学领域杰出贡献的正式肯定。1932年的这次大会常被认为是她职业生涯的巅峰。\(^\text{[129][131]}\)
\subsubsection{被纳粹德国驱逐出哥廷根大学}
1933年1月，阿道夫·希特勒出任德国总理后，纳粹活动在全国范围内急剧升级。在哥廷根大学，德国学生联合会带头攻击被他们称为“非德意志精神”的犹太学术传统，而这一行动还得到了诺特的前学生、当时的私人讲师维尔纳·韦伯的协助。反犹情绪弥漫，使得犹太教授在大学中陷入极端敌对的环境。据报道，一名年轻的抗议者甚至高喊：“雅利安学生要学习雅利安数学，而不是犹太数学！”\(^\text{[91]}\)

希特勒政府上台后的首批措施之一是《恢复职业公务员法》，该法令规定：凡是犹太人或政治上“不可靠”的公务人员（包括大学教授），除非在第一次世界大战中“以行动证明其对德国的忠诚”，否则必须被解职。1933年4月，普鲁士科学、艺术与公共教育部向诺特发出正式通知，内容为：“根据1933年4月7日《公务员法》第三条之规定，特此撤销您在哥廷根大学的授课权。”\(^\text{[132][133]}\)与诺特同时被解职的还有她的多位同事，包括马克斯·玻恩和理查德·库朗。\(^\text{[132][133]}\)

面对这场政治迫害，诺特表现得异常平静，她反而在困难时期向他人提供支持。赫尔曼·外尔后来写道：“艾米·诺特——她的勇气、坦率、对自身命运的淡然、以及她的宽容精神——在那充满仇恨、卑劣、绝望与悲伤的时代，成为了我们的精神慰藉。”\(^\text{[91]}\)如往常一样，诺特依然专注于数学研究。她在自己公寓中召集学生，继续讨论类域论。有一次，一名学生身穿纳粹准军事组织“冲锋队”的制服出现在聚会上，诺特并未表现出任何不安，甚至据说事后还对此一笑置之。\(^\text{[132][133]}\)
\subsubsection{避难于布林茅尔学院与普林斯顿}
当大批失业的德国教授开始在国外寻求职位时，他们在美国的同行纷纷伸出援手，为他们提供帮助与工作机会。阿尔伯特·爱因斯坦与赫尔曼·外尔被普林斯顿高等研究院聘用，而其他人则努力为流亡学者寻找合法移民所需的资助人。诺特收到了两所高校的邀请：一所是位于美国的布林茅尔学院，另一所是英国牛津大学的萨默维尔学院。经过与洛克菲勒基金会的一系列协商，基金会批准向布林茅尔学院提供资助，使诺特得以在那里任职，并于1933年底开始工作。\(^\text{[134][135]}\)

在布林茅尔学院，诺特结识并成为了安娜·惠勒的朋友，后者曾在诺特抵达前不久于哥廷根学习过。学院校长玛丽安·爱德华兹·帕克也是她的重要支持者，她热情地邀请周边学者前来“看看诺特博士的风采！”\(^\text{[136][137]}\)

在布林茅尔学院任职期间，诺特组建了一个学术小组，有时被称为“**诺特女孩**”（*Noether girls*）。[138] 该小组成员包括三位博士后——格蕾丝·肖弗·奎因（Grace Shover Quinn）、玛丽·约翰娜·魏斯（Marie Johanna Weiss）和奥尔加·陶斯基－托德（Olga Taussky-Todd），她们后来都在数学领域取得了杰出成就——以及一位博士生露丝·施托弗（Ruth Stauffer）。[139] 她们热情地研读了范德瓦尔登的《现代代数》（*Moderne Algebra I*）以及埃里希·黑克（Erich Hecke）的《代数数论》（*Theorie der algebraischen Zahlen*）的部分章节。[140]施托弗是诺特在美国唯一的博士生，但诺特在她毕业前不久去世。[141] 施托弗的论文答辩由理查德·布劳尔主持，并于1935年6月获得博士学位，[142] 论文主题为**可分正规扩张。[143] 获得博士学位后，施托弗曾短暂任教，随后在统计领域工作逾三十年。[46][142]

1934年，诺特受亚伯拉罕·弗莱克斯纳和奥斯瓦尔德·维布伦的邀请，开始在普林斯顿高等研究院讲学。[144] 她还与亚伯拉罕·阿尔伯特和哈里·范迪弗展开合作研究。[145] 对于普林斯顿大学，诺特曾半开玩笑地评论道，那是“一个男子的大学，任何女性都不得踏入”。[146]

诺特在美国的时光总体愉快，她被友善的同事所包围，能够全身心投入自己热爱的数学研究。[147] 1934年中期，她曾短暂返回德国，与埃米尔·阿廷及其弟弗里茨见面。[148] 弗里茨在被迫离开布雷斯劳工业大学后，接受了俄罗斯西伯利亚联邦区托木斯克数学与力学研究所的职位。[148]

尽管她的许多昔日同事都被逐出大学，诺特仍被允许以“外国学者”的身份使用哥廷根的图书馆。之后，她顺利返回美国，继续在布林茅尔学院进行教学与研究。[149][150]
